\documentclass[12pt]{article}
\def\activityid{F08-X-v3}\usepackage[letterpaper,margin=.65in,top=.60in,bottom=.65in]{geometry}
\usepackage[T1]{fontenc}
\usepackage[default]{sourcesanspro}
\usepackage{microtype,tikz,array,tabularx,fancyhdr,amsmath,amssymb}
\usepackage[hidelinks]{hyperref}
\usetikzlibrary{calc,arrows.meta}
\setlength{\parindent}{0pt}\setlength{\parskip}{7pt}
\setlength{\headheight}{0pt}\setlength{\footskip}{26pt}\setlength{\emergencystretch}{2em}
\pagestyle{fancy}\fancyhf{}\renewcommand{\headrulewidth}{0pt}\renewcommand{\footrulewidth}{.4pt}
\fancyfoot[L]{\scriptsize Bellingham Math Circle / Week 8 / \activityid}
\fancyfoot[R]{\scriptsize \thepage}
\definecolor{ink}{gray}{.12}\definecolor{soft}{gray}{.95}\color{ink}
\newcommand{\PageTitle}[2]{{\fontsize{10}{12}\selectfont\bfseries WEEK 8\quad STRATEGY LAB\hfill #1\par}\vspace{3pt}{\fontsize{26}{29}\selectfont\bfseries #2\par}\vspace{2pt}}
\newcommand{\Subhead}[1]{\par\vspace{3pt}{\large\bfseries #1}\par}
\newcommand{\RuleBox}[1]{\par\begingroup\setlength{\fboxsep}{8pt}\colorbox{soft}{\parbox{\dimexpr\linewidth-16pt\relax}{#1}}\endgroup\par}
\newcommand{\WriteBox}[2]{\par\textbf{#1}\par\vspace{2pt}\begin{tikzpicture}\draw[black!40,line width=.5pt] (0,0) rectangle (\dimexpr\linewidth-.5pt\relax,#2);\end{tikzpicture}\par}
\newcommand{\RookBoard}[3]{\begin{tikzpicture}[x=#3,y=#3]\draw[step=1,black!45] (0,0) grid (#1,#2);\draw[thick] (0,0) rectangle (#1,#2);\node[font=\Large] at ({#1-.5},.5){$\star$};\draw[-{Latex},thick] (0,-.35)--(#1,-.35);\draw[-{Latex},thick] ({#1+.35},#2)--({#1+.35},0);\end{tikzpicture}}
\newcommand{\Table}[3]{\begin{tikzpicture}[x=#3,y=#3]\draw[step=1,black!25] (0,0) grid (#1,#2);\draw[thick] (0,0) rectangle (#1,#2);\fill (0,0)circle(3pt);\draw[-{Latex},very thick] (0,0)--(.7,.7);\node[below] at ({#1/2},-.1){width #1};\node[rotate=90,above] at (-.2,{#2/2}){height #2};\end{tikzpicture}}
\newcommand{\GraphNodes}{\coordinate(A)at(0,0);\coordinate(B)at(3,0);\coordinate(C)at(1.5,2.6);\coordinate(D)at(1.5,.95);}
\newcommand{\Kfour}[1]{\begin{tikzpicture}[scale=#1]\GraphNodes\draw[very thick](A)--(B)--(C)--(A)--(D)--(B) (D)--(C);\foreach \n in {A,B,C,D}{\node[circle,draw,fill=white,minimum size=22pt,font=\bfseries]at(\n){\n};}\end{tikzpicture}}
\newcommand{\House}[1]{\begin{tikzpicture}[scale=#1]\coordinate(A)at(0,0);\coordinate(B)at(3,0);\coordinate(C)at(3,2);\coordinate(D)at(0,2);\coordinate(E)at(1.5,3.3);\draw[very thick](A)--(B)--(C)--(D)--(A) (D)--(E)--(C);\foreach \n in {A,B,C,D,E}{\node[circle,draw,fill=white,minimum size=22pt,font=\bfseries]at(\n){\n};}\end{tikzpicture}}




\newcommand{\StudentPage}[1]{{\fontsize{10}{12}\selectfont\bfseries Week 8 / Strategy lab / #1\par}\vspace{10pt}}
\newcommand{\Problem}[2]{\par\noindent\textbf{Problem #1:} #2\par}
\newcommand{\Space}[1]{\par\begin{tikzpicture}\draw[black!35] (0,0) rectangle (\dimexpr\linewidth-.5pt\relax,#1);\end{tikzpicture}\par}
\newcommand{\Piles}[1]{\begin{center}\begin{tikzpicture}[x=1in,y=1in]\foreach \x in {0,...,#1}{\draw[black!50,rounded corners] ({\x*2.05},0) rectangle ({\x*2.05+1.8},1.3);}\end{tikzpicture}\end{center}}

\begin{document}

\StudentPage{Extra / Grades 6--7}
\Problem{1}{Use two piles. On a turn, remove one or more counters from one pile OR remove the same positive number from both. Take the last counter to win. Play each start twice, changing who starts. Record one first move and reply.}
\Piles{1}
\begin{center}\renewcommand{\arraystretch}{2.3}\begin{tabular}{c|p{2.3in}|p{2.3in}}Start&First move leaves&Reply leaves\\\hline(2,2)&&\\(1,2)&&\\(2,3)&&\\(3,5)&&\end{tabular}\end{center}
\Problem{2}{From (1,2), try every legal first move. There are four. Draw or list them, then find a reply that wins immediately in each case.}
\Space{1.5in}
\newpage
\StudentPage{Extra / Grades 6--7}
\Problem{3}{Make a map to find starts worth giving your opponent. T means its next player loses against perfect replies; W means that player can force a win. Start with T at (0,0), where play is over. Work in increasing a+b. Mark W when you can move to a T. Mark T only when all legal moves go to W positions. Pile order does not matter; sort the new sizes before locating a move in this table.}
\begin{center}\renewcommand{\arraystretch}{2.1}\begin{tabular}{|c|*{8}{>{\centering\arraybackslash}p{.48in}|}}\hline$a\,\backslash\,b$&0&1&2&3&4&5&6&7\\\hline0&T&&&&&&&\\\hline 1&--&&&&&&&\\\hline 2&--&--&&&&&&\\\hline 3&--&--&--&&&&&\\\hline 4&--&--&--&--&&&&\\\hline 5&--&--&--&--&--&&&\\\hline 6&--&--&--&--&--&--&&\\\hline 7&--&--&--&--&--&--&--&\\\hline \end{tabular}\end{center}
\Problem{4}{List the T pairs you found. For each of (2,4), (4,6), and (5,7), build the piles and find a legal move to one of your T pairs. Record the moves.}
\Space{1.4in}
\newpage
\StudentPage{Extra / Grades 6--7}
\Problem{5}{Build a second list. Begin with (0,0). In row n, choose the smallest nonnegative number not yet used in either column as $a_n$. Set $b_n=a_n+n$. Complete the table, crossing numbers off the strip as you use them. Compare its pairs with your game map.}
\begin{center}0\quad1\quad2\quad3\quad4\quad5\quad6\quad7\quad8\quad9\quad10\quad11\quad12\quad13\end{center}
\begin{center}\renewcommand{\arraystretch}{1.85}\begin{tabular}{c|*{6}{>{\centering\arraybackslash}p{.65in}}}n&0&1&2&3&4&5\\\hline$a_n$&0&1&&&&\\$b_n$&0&2&&&&\\$b_n-a_n$&0&1&&&&\end{tabular}\end{center}
\Problem{6}{Choose two pairs in your list. Try to move from the larger pair to the smaller by taking from one pile, then by taking equally from both. Repeat for two other pairs. Record what prevents each move, or record a move if you find one.}
\Space{1.0in}
\Problem{7}{A known theorem gives the losing pairs as $(\lfloor n\phi\rfloor,\lfloor n\phi^2\rfloor)$, where $\phi=(1+\sqrt5)/2\approx1.618034$. The brackets mean round down. Use a calculator to check n=1,2,3,4,5 against your list. Then find moves into the list from (5,8) and (8,12). Record the moves.}
\Space{.7in}

\end{document}
